\ruleblock{Assurdo ($\bot$)}{
    \begin{prooftree}
        \hypo{\bot} 
        \infer1[($\bot E$)]{\phi}
    \end{prooftree}
}
{La regola di eliminazione esprime il principio "ex falso quodlibet". 
Non esiste una regola di introduzione poiché, secondo la semantica BHK, 
una contraddizione non ha dimostrazione costruttiva.}

\ruleblock{Implicazione ($\to$)}{
    \begin{minipage}{0.45\textwidth}
        \centering
        \begin{prooftree}
            \hypo{[\phi]^i} 
            \infer1{\psi} 
            \infer1[($\to I^i$)]{\phi \to \psi}
        \end{prooftree}
    \end{minipage}%
    \hfill
    \begin{minipage}{0.45\textwidth}
        \centering
        \begin{prooftree}
            \hypo{\phi \to \psi} 
            \hypo{\phi} 
            \infer2[($\to E$)]{\psi}
        \end{prooftree}
    \end{minipage}
}{La regola di introduzione estingue il debito logico $[\phi]^i$ internalizzandolo in 
un costruttore al di sotto della riga. Viceversa, la regola di eliminazione consuma 
l'implicazione al di sopra della riga per rilasciare la conclusione.}

\ruleblock{Congiunzione ($\land$)}{
    \begin{minipage}{0.38\textwidth}
        \centering
        \begin{prooftree}
            \hypo{\phi} 
            \hypo{\psi} 
            \infer2[($\land I$)]{\phi \land \psi}
        \end{prooftree}
    \end{minipage}%
    \hfill
    \begin{minipage}{0.58\textwidth}
        \centering
        \begin{prooftree}
            \hypo{\phi \land \psi} 
            \infer1[($\land E_1$)]{\phi}
        \end{prooftree}
        \quad
        \begin{prooftree}
            \hypo{\phi \land \psi} 
            \infer1[($\land E_2$)]{\psi}
        \end{prooftree}
    \end{minipage}
}{La regola di introduzione sintetizza una coppia di verità al di sotto della linea.
Le regole di eliminazione agiscono come proiezioni, analizzando e decostruendo il
termine posto al di sopra della linea.}

\ruleblock{Disgiunzione ($\lor$)}{
    \begin{minipage}{0.48\textwidth}
        \centering
        \begin{prooftree}
            \hypo{\phi} 
            \infer1[($\lor I_1$)]{\phi \lor \psi}
        \end{prooftree}
        \qquad
        \begin{prooftree}
            \hypo{\psi} 
            \infer1[($\lor I_2$)]{\phi \lor \psi}
        \end{prooftree}
    \end{minipage}%
    \hfill
    \begin{minipage}{0.48\textwidth} 
        \centering
        \small
        \begin{prooftree}
            \hypo{\phi \lor \psi} 
            \hypo{[\phi]^i} \infer1{\chi} 
            \quad 
            \hypo{[\psi]^j} \infer1{\chi} 
            \infer3[($\lor E^{i,j}$)]{\chi}
        \end{prooftree}
    \end{minipage}
}{La regola di eliminazione esegue un'analisi dei casi sul costruttore sopra 
la riga, estinguendo simultaneamente due debiti ipotetici ($i, j$) per 
raggiungere una conclusione comune.}